% --- Archon physics formalization source begin ---
% archon:physics
% archon:covers PhyXMiniProblems/problem_phyx_mini_0780.lean
% archon:source-report reports/phyx_mini/problem_phyx_mini_0780.source.json

\chapter{Physics problem phyx\_mini\_0780}
\label{ch:PhyXMiniProblems_problem_phyx_mini_0780}

\paragraph{Problem source.}
\#\# Physical scenario

A thin, light string is wrapped around the outer rim of a uniform hollow cylinder of mass \textbackslash{}(4.75\textbackslash{}text\{ kg\}\textbackslash{}) having inner and outer radii as shown in figure. The cylinder is then released from rest.

\#\# Question

How far must the cylinder fall before its center is moving at \textbackslash{}(6.66\textbackslash{}text\{ m/s\}\textbackslash{})?

\#\# Answer choices

- A: A: 4.76m

- B: B: 3.76m

- C: C: 2.76m

- D: D: 1.76m

\#\# Figure caption (auxiliary; use the image as primary evidence)

The image depicts a suspended orange ring. It consists of:

1. **Ring Structure**: An orange circular ring with a visible central hole.
   
2. **Suspension**: A thin gray cable or string attached to the upper edge of the ring, connected to a structure above (a horizontal bar).

3. **Measurements**:
   - An inner radius marked as 20.0 cm.
   - An outer radius marked as 35.0 cm.
   
4. **Arrows and Labels**: Two arrows pointing from the center of the ring outward to indicate the measurements, with text labeling the distances.

Overall, it illustrates a hanging ring with specified inner and outer radii.

\#\# Dataset metadata

Rotational Motion; Physical Model Grounding Reasoning, Multi-Formula Reasoning

\paragraph{Recorded answer/context.}
B: B: 3.76m

\paragraph{Figure/image path.}
/root/proposal\_for\_physic/science-mango/phyx\_mini\_run/phyx\_data/test\_image/780.png

\paragraph{Formalization target.}
create a compiling Lean file with sorry bodies at `PhyXMiniProblems/problem\_phyx\_mini\_0780.lean`. The Lean declarations must preserve the physical quantities, dimensions or dimensional roles, figure labels, governing-law hypotheses, and final relation expressed by this problem.
Use Mathlib/Physlib names found through LeanExplore where available. If a physics API is missing, introduce faithful local abstractions rather than scalar placeholder aliases.

\begin{theorem}[Physics formalization target]
\label{thm:physics:phyx_mini_0780:target}
\lean{PhyXMiniProblems.ProblemPhyXMini0780.distanceAtTargetSpeed_matches_answerB}
\uses{def:physics:phyx-mini-0780:phyxminiproblems-problemphyxmini0780-fallinghollowcylindersetup, def:physics:phyx-mini-0780:phyxminiproblems-problemphyxmini0780-matchesprimaryfigure, def:physics:phyx-mini-0780:phyxminiproblems-problemphyxmini0780-matchesproblemdata, def:physics:phyx-mini-0780:phyxminiproblems-problemphyxmini0780-hasphysicalparameters, def:physics:phyx-mini-0780:phyxminiproblems-problemphyxmini0780-satisfiesfallinghollowcylinderlaws, def:physics:phyx-mini-0780:phyxminiproblems-problemphyxmini0780-matchesdisplayeddistance}
The assigned autoformalize agent should translate this physics problem into Lean declarations in the covered file, with theorem and lemma proof bodies written as `by sorry`.
\end{theorem}
\begin{proof}
This is an autoformalization task, not a proof task. Produce statements that can later be proved by the physics prover without weakening the physical model.
\end{proof}

% --- Archon declaration topology begin ---
\section{Declaration topology}
\begin{definition}[Mass Quantity]
  \label{def:physics:phyx-mini-0780:phyxminiproblems-problemphyxmini0780-massquantity}
  \lean{PhyXMiniProblems.ProblemPhyXMini0780.MassQuantity}
  The nonnegative physical mass of the hollow cylinder
\end{definition}
\begin{proof}
  This is the declared model definition of Mass Quantity.
\end{proof}
\begin{definition}[Length Quantity]
  \label{def:physics:phyx-mini-0780:phyxminiproblems-problemphyxmini0780-lengthquantity}
  \lean{PhyXMiniProblems.ProblemPhyXMini0780.LengthQuantity}
  A nonnegative length magnitude, used for both radii and fall distance
\end{definition}
\begin{proof}
  This is the declared model definition of Length Quantity.
\end{proof}
\begin{definition}[Speed Quantity]
  \label{def:physics:phyx-mini-0780:phyxminiproblems-problemphyxmini0780-speedquantity}
  \lean{PhyXMiniProblems.ProblemPhyXMini0780.SpeedQuantity}
  A nonnegative center-of-mass speed magnitude
\end{definition}
\begin{proof}
  This is the declared model definition of Speed Quantity.
\end{proof}
\begin{definition}[Angular Speed Quantity]
  \label{def:physics:phyx-mini-0780:phyxminiproblems-problemphyxmini0780-angularspeedquantity}
  \lean{PhyXMiniProblems.ProblemPhyXMini0780.AngularSpeedQuantity}
  An angular-speed magnitude, with physical dimension inverse time
\end{definition}
\begin{proof}
  This is the declared model definition of Angular Speed Quantity.
\end{proof}
\begin{definition}[Acceleration Quantity]
  \label{def:physics:phyx-mini-0780:phyxminiproblems-problemphyxmini0780-accelerationquantity}
  \lean{PhyXMiniProblems.ProblemPhyXMini0780.AccelerationQuantity}
  An acceleration magnitude, with physical dimension L T⁻²
\end{definition}
\begin{proof}
  This is the declared model definition of Acceleration Quantity.
\end{proof}
\begin{definition}[Moment Of Inertia Quantity]
  \label{def:physics:phyx-mini-0780:phyxminiproblems-problemphyxmini0780-momentofinertiaquantity}
  \lean{PhyXMiniProblems.ProblemPhyXMini0780.MomentOfInertiaQuantity}
  A scalar moment of inertia about the cylinder's symmetry axis
\end{definition}
\begin{proof}
  This is the declared model definition of Moment Of Inertia Quantity.
\end{proof}
\begin{definition}[Energy Quantity]
  \label{def:physics:phyx-mini-0780:phyxminiproblems-problemphyxmini0780-energyquantity}
  \lean{PhyXMiniProblems.ProblemPhyXMini0780.EnergyQuantity}
  A real-valued physical energy
\end{definition}
\begin{proof}
  This is the declared model definition of Energy Quantity.
\end{proof}
\begin{definition}[mass In Kilograms]
  \label{def:physics:phyx-mini-0780:phyxminiproblems-problemphyxmini0780-massinkilograms}
  \lean{PhyXMiniProblems.ProblemPhyXMini0780.massInKilograms}
  \uses{def:physics:phyx-mini-0780:phyxminiproblems-problemphyxmini0780-massquantity}
  Kilogram readout of a physical mass
\end{definition}
\begin{proof}
  This is the declared model definition of mass In Kilograms.
\end{proof}
\begin{definition}[length In Meters]
  \label{def:physics:phyx-mini-0780:phyxminiproblems-problemphyxmini0780-lengthinmeters}
  \lean{PhyXMiniProblems.ProblemPhyXMini0780.lengthInMeters}
  \uses{def:physics:phyx-mini-0780:phyxminiproblems-problemphyxmini0780-lengthquantity}
  Metre readout of a physical length
\end{definition}
\begin{proof}
  This is the declared model definition of length In Meters.
\end{proof}
\begin{definition}[speed In Meters Per Second]
  \label{def:physics:phyx-mini-0780:phyxminiproblems-problemphyxmini0780-speedinmeterspersecond}
  \lean{PhyXMiniProblems.ProblemPhyXMini0780.speedInMetersPerSecond}
  \uses{def:physics:phyx-mini-0780:phyxminiproblems-problemphyxmini0780-speedquantity}
  Metre-per-second readout of a physical speed
\end{definition}
\begin{proof}
  This is the declared model definition of speed In Meters Per Second.
\end{proof}
\begin{definition}[angular Speed In Radians Per Second]
  \label{def:physics:phyx-mini-0780:phyxminiproblems-problemphyxmini0780-angularspeedinradianspersecond}
  \lean{PhyXMiniProblems.ProblemPhyXMini0780.angularSpeedInRadiansPerSecond}
  \uses{def:physics:phyx-mini-0780:phyxminiproblems-problemphyxmini0780-angularspeedquantity}
  Radian-per-second readout of an angular-speed magnitude
\end{definition}
\begin{proof}
  This is the declared model definition of angular Speed In Radians Per Second.
\end{proof}
\begin{definition}[acceleration In Meters Per Second Squared]
  \label{def:physics:phyx-mini-0780:phyxminiproblems-problemphyxmini0780-accelerationinmeterspersecondsquared}
  \lean{PhyXMiniProblems.ProblemPhyXMini0780.accelerationInMetersPerSecondSquared}
  \uses{def:physics:phyx-mini-0780:phyxminiproblems-problemphyxmini0780-accelerationquantity}
  Metre-per-second-squared readout of an acceleration magnitude
\end{definition}
\begin{proof}
  This is the declared model definition of acceleration In Meters Per Second Squared.
\end{proof}
\begin{definition}[moment Of Inertia In Kilogram Meters Squared]
  \label{def:physics:phyx-mini-0780:phyxminiproblems-problemphyxmini0780-momentofinertiainkilogrammeterssquared}
  \lean{PhyXMiniProblems.ProblemPhyXMini0780.momentOfInertiaInKilogramMetersSquared}
  \uses{def:physics:phyx-mini-0780:phyxminiproblems-problemphyxmini0780-momentofinertiaquantity}
  Kilogram-metre-squared readout of a scalar moment of inertia
\end{definition}
\begin{proof}
  This is the declared model definition of moment Of Inertia In Kilogram Meters Squared.
\end{proof}
\begin{definition}[energy In Joules]
  \label{def:physics:phyx-mini-0780:phyxminiproblems-problemphyxmini0780-energyinjoules}
  \lean{PhyXMiniProblems.ProblemPhyXMini0780.energyInJoules}
  \uses{def:physics:phyx-mini-0780:phyxminiproblems-problemphyxmini0780-energyquantity}
  Joule readout of a physical energy
\end{definition}
\begin{proof}
  This is the declared model definition of energy In Joules.
\end{proof}
\begin{definition}[Motion Event]
  \label{def:physics:phyx-mini-0780:phyxminiproblems-problemphyxmini0780-motionevent}
  \lean{PhyXMiniProblems.ProblemPhyXMini0780.MotionEvent}
  The release event and the later event specified by the target speed
\end{definition}
\begin{proof}
  This is the declared model definition of Motion Event.
\end{proof}
\begin{definition}[Figure Object]
  \label{def:physics:phyx-mini-0780:phyxminiproblems-problemphyxmini0780-figureobject}
  \lean{PhyXMiniProblems.ProblemPhyXMini0780.FigureObject}
  Physical objects visible in the supplied bitmap
\end{definition}
\begin{proof}
  This is the declared model definition of Figure Object.
\end{proof}
\begin{definition}[Figure Label]
  \label{def:physics:phyx-mini-0780:phyxminiproblems-problemphyxmini0780-figurelabel}
  \lean{PhyXMiniProblems.ProblemPhyXMini0780.FigureLabel}
  The two numerical radius labels printed in the supplied bitmap
\end{definition}
\begin{proof}
  This is the declared model definition of Figure Label.
\end{proof}
\begin{definition}[Radius Kind]
  \label{def:physics:phyx-mini-0780:phyxminiproblems-problemphyxmini0780-radiuskind}
  \lean{PhyXMiniProblems.ProblemPhyXMini0780.RadiusKind}
  Which cylindrical surface a radius label identifies
\end{definition}
\begin{proof}
  This is the declared model definition of Radius Kind.
\end{proof}
\begin{definition}[Orientation]
  \label{def:physics:phyx-mini-0780:phyxminiproblems-problemphyxmini0780-orientation}
  \lean{PhyXMiniProblems.ProblemPhyXMini0780.Orientation}
  Orientations needed to transcribe the string and support
\end{definition}
\begin{proof}
  This is the declared model definition of Orientation.
\end{proof}
\begin{definition}[Hollow Cylinder Figure]
  \label{def:physics:phyx-mini-0780:phyxminiproblems-problemphyxmini0780-hollowcylinderfigure}
  \lean{PhyXMiniProblems.ProblemPhyXMini0780.HollowCylinderFigure}
  \uses{def:physics:phyx-mini-0780:phyxminiproblems-problemphyxmini0780-figureobject, def:physics:phyx-mini-0780:phyxminiproblems-problemphyxmini0780-figurelabel, def:physics:phyx-mini-0780:phyxminiproblems-problemphyxmini0780-radiuskind, def:physics:phyx-mini-0780:phyxminiproblems-problemphyxmini0780-orientation}
  Direct qualitative and label evidence from the primary bitmap
\end{definition}
\begin{proof}
  This is the declared model definition of Hollow Cylinder Figure.
\end{proof}
\begin{definition}[Cylinder Mass Distribution]
  \label{def:physics:phyx-mini-0780:phyxminiproblems-problemphyxmini0780-cylindermassdistribution}
  \lean{PhyXMiniProblems.ProblemPhyXMini0780.CylinderMassDistribution}
  The cylinder's mass distribution stated in the problem
\end{definition}
\begin{proof}
  This is the declared model definition of Cylinder Mass Distribution.
\end{proof}
\begin{definition}[String Mass Model]
  \label{def:physics:phyx-mini-0780:phyxminiproblems-problemphyxmini0780-stringmassmodel}
  \lean{PhyXMiniProblems.ProblemPhyXMini0780.StringMassModel}
  Idealized string-mass models
\end{definition}
\begin{proof}
  This is the declared model definition of String Mass Model.
\end{proof}
\begin{definition}[String Thickness Model]
  \label{def:physics:phyx-mini-0780:phyxminiproblems-problemphyxmini0780-stringthicknessmodel}
  \lean{PhyXMiniProblems.ProblemPhyXMini0780.StringThicknessModel}
  Idealized transverse-thickness models for the string
\end{definition}
\begin{proof}
  This is the declared model definition of String Thickness Model.
\end{proof}
\begin{definition}[String Contact Regime]
  \label{def:physics:phyx-mini-0780:phyxminiproblems-problemphyxmini0780-stringcontactregime}
  \lean{PhyXMiniProblems.ProblemPhyXMini0780.StringContactRegime}
  How the string and cylinder move relative to one another
\end{definition}
\begin{proof}
  This is the declared model definition of String Contact Regime.
\end{proof}
\begin{definition}[Release Protocol]
  \label{def:physics:phyx-mini-0780:phyxminiproblems-problemphyxmini0780-releaseprotocol}
  \lean{PhyXMiniProblems.ProblemPhyXMini0780.ReleaseProtocol}
  How the cylinder begins its motion
\end{definition}
\begin{proof}
  This is the declared model definition of Release Protocol.
\end{proof}
\begin{definition}[Falling Hollow Cylinder Setup]
  \label{def:physics:phyx-mini-0780:phyxminiproblems-problemphyxmini0780-fallinghollowcylindersetup}
  \lean{PhyXMiniProblems.ProblemPhyXMini0780.FallingHollowCylinderSetup}
  \uses{def:physics:phyx-mini-0780:phyxminiproblems-problemphyxmini0780-massquantity, def:physics:phyx-mini-0780:phyxminiproblems-problemphyxmini0780-lengthquantity, def:physics:phyx-mini-0780:phyxminiproblems-problemphyxmini0780-speedquantity, def:physics:phyx-mini-0780:phyxminiproblems-problemphyxmini0780-angularspeedquantity, def:physics:phyx-mini-0780:phyxminiproblems-problemphyxmini0780-accelerationquantity, def:physics:phyx-mini-0780:phyxminiproblems-problemphyxmini0780-momentofinertiaquantity, def:physics:phyx-mini-0780:phyxminiproblems-problemphyxmini0780-energyquantity, def:physics:phyx-mini-0780:phyxminiproblems-problemphyxmini0780-motionevent, def:physics:phyx-mini-0780:phyxminiproblems-problemphyxmini0780-hollowcylinderfigure, def:physics:phyx-mini-0780:phyxminiproblems-problemphyxmini0780-cylindermassdistribution, def:physics:phyx-mini-0780:phyxminiproblems-problemphyxmini0780-stringmassmodel, def:physics:phyx-mini-0780:phyxminiproblems-problemphyxmini0780-stringthicknessmodel, def:physics:phyx-mini-0780:phyxminiproblems-problemphyxmini0780-stringcontactregime, def:physics:phyx-mini-0780:phyxminiproblems-problemphyxmini0780-releaseprotocol}
  The physical quantities and SI rigid-body data used to model the two events. The fall distance is an independent physical quantity: it is not defined from the requested 3.76 m answer
\end{definition}
\begin{proof}
  This is the declared model definition of Falling Hollow Cylinder Setup.
\end{proof}
\begin{definition}[Matches Primary Figure]
  \label{def:physics:phyx-mini-0780:phyxminiproblems-problemphyxmini0780-matchesprimaryfigure}
  \lean{PhyXMiniProblems.ProblemPhyXMini0780.MatchesPrimaryFigure}
  \uses{def:physics:phyx-mini-0780:phyxminiproblems-problemphyxmini0780-figureobject, def:physics:phyx-mini-0780:phyxminiproblems-problemphyxmini0780-figurelabel, def:physics:phyx-mini-0780:phyxminiproblems-problemphyxmini0780-fallinghollowcylindersetup}
  Primary-image evidence: the cylinder is annular, the string is vertical and tangent to the outer rim, its upper end meets a horizontal support, and arrows from the center label the inner and outer radii as 20.0 cm and 35.0 cm
\end{definition}
\begin{proof}
  This is the declared model definition of Matches Primary Figure.
\end{proof}
\begin{definition}[Matches Problem Data]
  \label{def:physics:phyx-mini-0780:phyxminiproblems-problemphyxmini0780-matchesproblemdata}
  \lean{PhyXMiniProblems.ProblemPhyXMini0780.MatchesProblemData}
  \uses{def:physics:phyx-mini-0780:phyxminiproblems-problemphyxmini0780-massinkilograms, def:physics:phyx-mini-0780:phyxminiproblems-problemphyxmini0780-lengthinmeters, def:physics:phyx-mini-0780:phyxminiproblems-problemphyxmini0780-speedinmeterspersecond, def:physics:phyx-mini-0780:phyxminiproblems-problemphyxmini0780-angularspeedinradianspersecond, def:physics:phyx-mini-0780:phyxminiproblems-problemphyxmini0780-accelerationinmeterspersecondsquared, def:physics:phyx-mini-0780:phyxminiproblems-problemphyxmini0780-fallinghollowcylindersetup}
  Numerical and qualitative data supplied by the statement and primary figure. The radius readouts convert 20.0 cm and 35.0 cm to 1/5 m and 7/20 m. The standard terrestrial acceleration readout 9.81 m/s² is made explicit so that the rounding convention behind the multiple-choice answer is visible
\end{definition}
\begin{proof}
  This is the declared model definition of Matches Problem Data.
\end{proof}
\begin{definition}[Has Physical Parameters]
  \label{def:physics:phyx-mini-0780:phyxminiproblems-problemphyxmini0780-hasphysicalparameters}
  \lean{PhyXMiniProblems.ProblemPhyXMini0780.HasPhysicalParameters}
  \uses{def:physics:phyx-mini-0780:phyxminiproblems-problemphyxmini0780-massinkilograms, def:physics:phyx-mini-0780:phyxminiproblems-problemphyxmini0780-lengthinmeters, def:physics:phyx-mini-0780:phyxminiproblems-problemphyxmini0780-accelerationinmeterspersecondsquared, def:physics:phyx-mini-0780:phyxminiproblems-problemphyxmini0780-fallinghollowcylindersetup}
  Positivity and nondegeneracy of the independent physical parameters
\end{definition}
\begin{proof}
  This is the declared model definition of Has Physical Parameters.
\end{proof}
\begin{definition}[Satisfies Falling Hollow Cylinder Laws]
  \label{def:physics:phyx-mini-0780:phyxminiproblems-problemphyxmini0780-satisfiesfallinghollowcylinderlaws}
  \lean{PhyXMiniProblems.ProblemPhyXMini0780.SatisfiesFallingHollowCylinderLaws}
  \uses{def:physics:phyx-mini-0780:phyxminiproblems-problemphyxmini0780-massinkilograms, def:physics:phyx-mini-0780:phyxminiproblems-problemphyxmini0780-lengthinmeters, def:physics:phyx-mini-0780:phyxminiproblems-problemphyxmini0780-speedinmeterspersecond, def:physics:phyx-mini-0780:phyxminiproblems-problemphyxmini0780-angularspeedinradianspersecond, def:physics:phyx-mini-0780:phyxminiproblems-problemphyxmini0780-accelerationinmeterspersecondsquared, def:physics:phyx-mini-0780:phyxminiproblems-problemphyxmini0780-momentofinertiainkilogrammeterssquared, def:physics:phyx-mini-0780:phyxminiproblems-problemphyxmini0780-energyinjoules, def:physics:phyx-mini-0780:phyxminiproblems-problemphyxmini0780-motionevent, def:physics:phyx-mini-0780:phyxminiproblems-problemphyxmini0780-fallinghollowcylindersetup}
  The governing rigid-body, kinematic, and energy laws. A uniform hollow cylinder has axial inertia I = (1/2) M (R\_inner² + R\_outer²). The thin string unwinds without slip, so v = R\_outer ω. Translational kinetic energy is M v² / 2, while rotational kinetic energy is linked to Physlib's inertia-tensor definition.
\end{definition}
\begin{proof}
  This is the declared model definition of Satisfies Falling Hollow Cylinder Laws.
\end{proof}
\begin{lemma}[distance At Target Speed exact]
  \label{lem:physics:phyx-mini-0780:phyxminiproblems-problemphyxmini0780-distanceattargetspeed-exact}
  \lean{PhyXMiniProblems.ProblemPhyXMini0780.distanceAtTargetSpeed_exact}
  \uses{def:physics:phyx-mini-0780:phyxminiproblems-problemphyxmini0780-lengthinmeters, def:physics:phyx-mini-0780:phyxminiproblems-problemphyxmini0780-fallinghollowcylindersetup, def:physics:phyx-mini-0780:phyxminiproblems-problemphyxmini0780-matchesproblemdata, def:physics:phyx-mini-0780:phyxminiproblems-problemphyxmini0780-hasphysicalparameters, def:physics:phyx-mini-0780:phyxminiproblems-problemphyxmini0780-satisfiesfallinghollowcylinderlaws}
  With g = 9.81 m/s², the exact ideal-model distance is 2008323 / 534100 m, approximately 3.760200337 m
\end{lemma}
\begin{proof}
  The conclusion follows from the governing relations and figure readouts named in the statement of distance At Target Speed exact.
\end{proof}
\begin{definition}[Rounds To Nearest Hundredth Meter]
  \label{def:physics:phyx-mini-0780:phyxminiproblems-problemphyxmini0780-roundstonearesthundredthmeter}
  \lean{PhyXMiniProblems.ProblemPhyXMini0780.RoundsToNearestHundredthMeter}
  \uses{def:physics:phyx-mini-0780:phyxminiproblems-problemphyxmini0780-lengthquantity, def:physics:phyx-mini-0780:phyxminiproblems-problemphyxmini0780-lengthinmeters}
  A displayed metre value is a valid nearest-hundredth approximation
\end{definition}
\begin{proof}
  This is the declared model definition of Rounds To Nearest Hundredth Meter.
\end{proof}
\begin{definition}[Answer Choice]
  \label{def:physics:phyx-mini-0780:phyxminiproblems-problemphyxmini0780-answerchoice}
  \lean{PhyXMiniProblems.ProblemPhyXMini0780.AnswerChoice}
  Labels of the four displayed answer choices
\end{definition}
\begin{proof}
  This is the declared model definition of Answer Choice.
\end{proof}
\begin{definition}[meters]
  \label{def:physics:phyx-mini-0780:phyxminiproblems-problemphyxmini0780-answerchoice-meters}
  \lean{PhyXMiniProblems.ProblemPhyXMini0780.AnswerChoice.meters}
  \uses{def:physics:phyx-mini-0780:phyxminiproblems-problemphyxmini0780-answerchoice}
  Distance in metres printed beside each displayed answer label
\end{definition}
\begin{proof}
  This is the declared model definition of meters.
\end{proof}
\begin{definition}[recorded Answer Choice]
  \label{def:physics:phyx-mini-0780:phyxminiproblems-problemphyxmini0780-recordedanswerchoice}
  \lean{PhyXMiniProblems.ProblemPhyXMini0780.recordedAnswerChoice}
  \uses{def:physics:phyx-mini-0780:phyxminiproblems-problemphyxmini0780-answerchoice}
  Dataset metadata recording the supplied answer label
\end{definition}
\begin{proof}
  This is the declared model definition of recorded Answer Choice.
\end{proof}
\begin{definition}[Matches Displayed Distance]
  \label{def:physics:phyx-mini-0780:phyxminiproblems-problemphyxmini0780-matchesdisplayeddistance}
  \lean{PhyXMiniProblems.ProblemPhyXMini0780.MatchesDisplayedDistance}
  \uses{def:physics:phyx-mini-0780:phyxminiproblems-problemphyxmini0780-lengthquantity, def:physics:phyx-mini-0780:phyxminiproblems-problemphyxmini0780-roundstonearesthundredthmeter, def:physics:phyx-mini-0780:phyxminiproblems-problemphyxmini0780-answerchoice}
  A physical distance rounds to the metre value displayed for a choice
\end{definition}
\begin{proof}
  This is the declared model definition of Matches Displayed Distance.
\end{proof}
\begin{theorem}[distance At Target Speed rounds To three Point Seven Six Meters]
  \label{thm:physics:phyx-mini-0780:phyxminiproblems-problemphyxmini0780-distanceattargetspeed-roundsto-threepointsevensixmeters}
  \lean{PhyXMiniProblems.ProblemPhyXMini0780.distanceAtTargetSpeed_roundsTo_threePointSevenSixMeters}
  \uses{def:physics:phyx-mini-0780:phyxminiproblems-problemphyxmini0780-fallinghollowcylindersetup, def:physics:phyx-mini-0780:phyxminiproblems-problemphyxmini0780-matchesprimaryfigure, def:physics:phyx-mini-0780:phyxminiproblems-problemphyxmini0780-matchesproblemdata, def:physics:phyx-mini-0780:phyxminiproblems-problemphyxmini0780-hasphysicalparameters, def:physics:phyx-mini-0780:phyxminiproblems-problemphyxmini0780-satisfiesfallinghollowcylinderlaws, def:physics:phyx-mini-0780:phyxminiproblems-problemphyxmini0780-roundstonearesthundredthmeter}
  The cylinder must fall about 3.76 m before its center reaches 6.66 m/s. This is answer choice B and formalizes
\end{theorem}
\begin{proof}
  The conclusion follows from the governing relations and figure readouts named in the statement of distance At Target Speed rounds To three Point Seven Six Meters.
\end{proof}
% --- Archon declaration topology end ---
% --- Archon physics formalization source end ---
